\chapter*{Abstract}
\thispagestyle{specialheader}
\markboth{}{}
\addcontentsline{toc}{chapter}{Abstract}
This thesis presents a multi-objective optimisation of a compact 48 MWth
pressurised water reactor core fuelled with low-enriched \ce{UO2}, based
on surrogate models and the OpenMC Monte Carlo code. The objective was to find
core configurations that meet the publicly reported requirements of the land-based prototype of the Brazilian naval propulsion
programme, and to quantify the trade-off
between cycle length, radial peaking and critical boron concentration.

The core data of the prototype are not public, so the model was built from the reactor type, the power, the fuel class, the vessel envelope and the five-year refuelling interval reported in open sources, and from declared assumptions. The enrichment, the gadolinia loading, the number of poisoned pins and the reflector
thickness were varied on a single-batch core of 32 assemblies. Because one
depletion evaluation takes around 20 minutes and has a statistical
error, a Gaussian-process surrogate model trained on the Monte Carlo
evaluations was searched by the NSGA-II multi-objective genetic algorithm
in an active-learning loop over nine campaigns and
636 evaluations.

Campaigns 1 to 8 maximised the cycle length and minimised the radial
peaking factor. Every design of the Campaign 8 front satisfied the five-year requirement on the assembly depletion, and the longer cycles required excess reactivity that only the control banks at a large insertion depth could compensate. Campaign 9 therefore imposed the cycle length as a constraint at 1826 effective full power days and minimised the peaking factor and the critical boron concentration at the most reactive state.

The five members of the Campaign 9 front satisfy every constraint of the search on the assembly depletion. In the post-analysis, the eight-layer depletion of the 3D core shortens every cycle by 15 to 25\,\%, and no member of that front remains feasible. The search was therefore continued on the axially corrected cycle length and converged to a new front, C9-27, C9-69 and C9-70, with cycle margins of 24 to 100 days, 3D peaking factors of 1.496 to 1.565 and critical boron concentrations of 1633 to 2268 ppm, below the moderator temperature coefficient limit of each lattice. All three are kept subcritical by the first two regulating banks, and C9-27 is the recommended design. These results assume fixed temperatures and constant boron in the depletion.
